\begin{corollary}
The set of formulas of a language with a countable alphabet is countable.
\end{corollary}
\begin{proof}
Every formula is a finite string over the alphabet, so the formula set is a
subset of the countable set \(A^*\).
\end{proof}

If \(u\in A^m\) and \(v\in A^n\), their concatenation \(uv\in A^{m+n}\) is
the sequence defined by
\[
  (uv)(i)=
  \begin{cases}
    u(i),&i<m,\\
    v(i-m),&m\leq i<m+n.
  \end{cases}
\]
Concatenation is associative and \(\epsilon\) is its identity.

\begin{lemma}
For all strings \(u,v,w\in A^*\),
\((uv)w=u(vw)\), \(\epsilon u=u\), and \(u\epsilon=u\).
\end{lemma}
\begin{proof}
Two finite sequences are equal iff they have the same length and the same
value at every argument.  Write \(m=\length(u)\), \(n=\length(v)\), and
\(p=\length(w)\).

For associativity, first compare lengths.  By the definition of
concatenation, \(\length(uv)=m+n\), so \(\length((uv)w)=(m+n)+p\); likewise
\(\length(vw)=n+p\), so \(\length(u(vw))=m+(n+p)\).  Addition on \(\N\) is
associative, so both strings are functions on \(\set{0,\ldots,m+n+p-1}\).

Now fix an argument \(i<m+n+p\) and distinguish the three exhaustive and
pairwise disjoint intervals \(i<m\), \(m\leq i<m+n\), and
\(m+n\leq i<m+n+p\).

If \(i<m\), then also \(i<m+n\), so
\(((uv)w)(i)=(uv)(i)=u(i)\); and \((u(vw))(i)=u(i)\) directly, since
\(i<m\).

If \(m\leq i<m+n\), then \(((uv)w)(i)=(uv)(i)=v(i-m)\), by the second
clause for \(uv\) applied at \(i\).  On the other side, \(i\geq m\) gives
\((u(vw))(i)=(vw)(i-m)\), and \(0\leq i-m<n\) puts \(i-m\) in the first
interval for \(vw\), so \((vw)(i-m)=v(i-m)\).

If \(m+n\leq i<m+n+p\), then \(i\geq m+n\) gives
\(((uv)w)(i)=w(i-(m+n))\).  On the other side, \(i\geq m\) gives
\((u(vw))(i)=(vw)(i-m)\), and \(i-m\geq n\) puts \(i-m\) in the second
interval for \(vw\), so \((vw)(i-m)=w((i-m)-n)=w(i-(m+n))\).

In each case the two strings agree at \(i\), so \((uv)w=u(vw)\).

For the identity equations, \(\length(\epsilon u)=0+m=m=\length(u)\), and
every argument \(i<m\) satisfies \(0\leq i<0+m\), so the second clause gives
\((\epsilon u)(i)=u(i-0)=u(i)\); hence \(\epsilon u=u\).  Similarly
\(\length(u\epsilon)=m+0=m\), and every argument \(i<m\) falls in the first
interval, so \((u\epsilon)(i)=u(i)\); hence \(u\epsilon=u\).
\end{proof}

